% Pursuing Stacks (cotes 134-2 to 134-8): archive cote 134-4, site page 5 (PDF page 6), typescript page 93.
% typescript. Draft compared: ps3.tex:205-276.
% AI-transcribed by Claude Opus 5.5 (claude-opus-5-5, adaptive thinking, effort high), Message Batches API, batch ps1 (msgbatch_01Rv5hCaNTtUt4JHx1yqoHPL), 2026-11-12; not yet checked by eye.
\scanpage{134-4/5}
\mspage{(93)}
46.   F\unsure{i}rst thing to do visibly is to define  the notion of a test-func\unsure{g|t}or
\[ i : A^{\unsure{*}} \longrightarrow M \quad , \]
where M is any modelizer. \xout{t}Thus M is endowed with a \xout{\unsure{famil}} subset $W_M < \mathrm{F}\unsure{l}(M)$
\xout{of} i?e. a notion of weak equivalence,  satisfying the ``mild saturation
conditions'' of p.59,  and moreover we assume that $W_M^{-1}M$  is equivalent to
(Hot)  -  but the choice of an equivalence,  or equivalently,  of the
corresponding localization functor
\[ M \longrightarrow (\mathrm{Hot}) \ , \]
is not given with the structure.  (If we admit the ``inspiring assumption'',
there is no real choice,  as a matter of fact  -  but we don't want to use
this in a technic\hand{\unsure{$\smile$}}al  sense,  but only as a guide and motivation.)

Let's start with the weak variant  -  we assume A to be a weak test
category,  and want to define what it means that i is a \uline{weak test functor}.
In all this game,  it is understood that in case M = (Cat),  the notions
we want to define (of a weak test functor and of a test functor) should
\unsure{b|r}educe to the ones we have pinpointed in yesterday's notes.

THe very first idea that comes to mind,  is to demand merely that
the corresponding \xout{it} $i^*$  (1)
\[ i^*:\ M \longrightarrow A^\wedge \]
should be  modelizing;$^{\unsure{\mathrm{x}}}$  which means (I recall)
\[ \text{a)}\quad W_M = i^{*\,-1}(\unsure{W_{A^\wedge}|W_A}) \]
\[ \text{b)}\quad \text{The induced functor } W_M^{-1}M \longrightarrow W_A^{-1}A^\wedge \text{ is an equivalence.} \]
This,  I just checked,  does correspond to the definition we gave yester-%
day (p.85),  when  M = (Cat).  THere is a very interesting extra feature
though in this special case,  which appears kind of ``in between the lines''
in the ``key result'' on p.61 , \inserted{and} which I want \inserted{now} to state in the more general
set up.

As usual in related situations,  the notion of weak equivalence in
M  gives rise to a corresponding notion of ``aspheric'' elements in M -
namely those for which the unique map
\[ x \longrightarrow e_M \]
is a weak equivalence. \xout{w}We assume now the existence of a final object $e_M$
in M ,  and will assume too,  if necessary,  that it's image in the
localization $W_M^{-1}M$ \inserted{$= H_M$} is equally a final object.\hand{$*$} \xout{If we assume} Thus \inserted{if} x in M
is aspheric,  it's image in $H_M$  is a final object,  and the converse holds
provided we assume $W_M$  strongly saturated\hand{,}namely any map in M which becomes
an isomorphism in $H_M$  is a weak equivalence.
\margin{$*$ 29.3. This condition will \inserted{be} verified, if there exists a weak test functor $i:A\to M$.}

\margin{5}
